% META: {"id": "TSPE-GEOESPACE-CR-010", "chapitre": "TSPE-GEOMETRIE-ESPACE", "type_objet": "cours", "capacites_codes": ["C1", "C2", "C4", "C5"], "sources_inspiration": [], "mode_creation": "ex_nihilo", "status": "approved"}
% BEGIN-VERIFY
% from sympy import Matrix
% A=Matrix([0,0,0]); B=Matrix([1,0,0]); D=Matrix([0,1,0]); E=Matrix([0,0,1])
% C=B+D-A; G=C+E-A
% assert G-A == (B-A)+(D-A)+(E-A)      # AG = AB + AD + AE
% assert C-B == D-A                    # BC = AD
% assert G-C == E-A                    # CG = AE
% assert list(G-A) == [1,1,1]          # coordonnees de AG dans la base
% END-VERIFY

\section{Vecteurs de l'espace, bases}

% ============================================================
% STRATE 1 — Combinaisons lineaires, bases
% ============================================================

\propriete[Combinaison linéaire]{
Etant donnes des vecteurs $\vec{u}, \vec{v}, \vec{w}$ de l'espace et des réels $a,b,c$, le vecteur $a\vec{u}+b\vec{v}+c\vec{w}$ est appele \textbf{combinaison linéaire} de $\vec{u},\vec{v},\vec{w}$. Les règles de calcul (addition de Chasles, multiplication par un réel) sont les mêmes que dans le plan.
}

\exemple{
Dans un cube $ABCDEFGH$, exprimer $\vec{AG}$ comme combinaison linéaire de $\vec{AB}$, $\vec{AD}$, $\vec{AE}$.

Par la relation de Chasles : $\vec{AG} = \vec{AB}+\vec{BC}+\vec{CG} = \vec{AB}+\vec{AD}+\vec{AE}$ (car $\vec{BC}=\vec{AD}$ et $\vec{CG}=\vec{AE}$ dans un cube).
}

\definition[1]{
Trois vecteurs $\vec{u},\vec{v},\vec{w}$ \textbf{non coplanaires} (c'est-a-dire tels qu'aucun n'est combinaison linéaire des deux autres) forment une \textbf{base} de l'espace. Tout vecteur $\vec{t}$ de l'espace se decompose alors de maniere \textbf{unique} :
\[
  \vec{t} = x\vec{u}+y\vec{v}+z\vec{w}, \qquad (x,y,z) \text{ appeles coordonnees de } \vec{t} \text{ dans la base } (\vec{u},\vec{v},\vec{w}).
\]
}

\margeAppui{
Sur une figure (par exemple un cube ou un parallelepipede), trois vecteurs non coplanaires forment une base si et seulement si ils ne sont pas tous les trois dans un même plan.
}

\exemple{
Dans le cube $ABCDEFGH$, $(\vec{AB},\vec{AD},\vec{AE})$ est une base de l'espace (trois arêtes issues du même sommet, non coplanaires). Les coordonnées de $\vec{AG}$ dans cette base sont $(1,1,1)$ (d'après l'exemple précédent).
}

% ============================================================
% STRATE 2 — Erreurs frequentes
% ============================================================

\erreurFrequente{
\textbf{Confondre coplanaire et colineaire.} Deux vecteurs sont toujours coplanaires (ils definissent, avec un point, un plan) ; la notion de coplanarite non-triviale concerne \emph{trois} vecteurs : ils sont coplanaires si l'un est combinaison linéaire des deux autres.
}
